% Part: computability
% Chapter: computability-theory
% Section: k-1

\documentclass[../../../include/open-logic-section]{subfiles}

\begin{document}

\olfileid{cmp}{thy}{k1}
\olsection{归约性的一个例子}

考虑 \olref[ppr]{prop:reduce} 的一个应用。

\begin{prop}
\ollabel{prop:k1}
设
\[
K_1 = \Setabs{e}{\cfind{e}(0) \fdefined}.
\]则 $K_1$ 是可计算枚举的，但不是可计算的。
\end{prop}

\begin{proof}
由于 $K_1 = \Setabs{e}{\lexists[s][T(e,0,s)]}$，根据 \olref[eqc]{thm:exists-char}，$K_1$ 是可计算枚举的。

为证 $K_1$ 不可计算，我们来证明 $K_0$ 可归约到 $K_1$。

\begin{explain}
这有一点技巧性，因为利用 $K_1$ 我们只能询问关于以特定输入 $0$ 开始的计算的问题。假设你有一个聪明的朋友，他能回答这类问题（这样的朋友被称为“预言机”）。此时，假设有人来问你 $\tuple{e, x}$ 是否属于 $K_0$，即机器 $e$ 在输入 $x$ 上是否停机。你可以构造另一台机器 $e_x$，对于\emph{任何}输入，它都忽略该输入，转而在输入 $x$ 上运行机器 $e$。于是，机器 $e$ 在输入 $x$ 上是否停机的问题，显然等价于机器 $e_x$ 在输入 $0$（或任何其他输入）上是否停机的问题。接着，你问你的朋友这台新机器 $e_x$ 在输入 $0$ 上是否停机；你朋友对修改后问题的回答，就给出了原问题的答案。这就实现了将 $K_0$ 归约到 $K_1$。
\end{explain}

利用通用部分可计算函数，令 $f$ 为如下定义的三元函数：
\[
f(x,y,z) \simeq \cfind{x}(y).
\]
注意 $f$ 完全忽略其第三个输入。选取一个指标 $e$ 使得 $f = \cfind{e}[3]$；于是我们有
\[
\cfind{e}[3](x,y,z) \simeq \cfind{x}(y).
\]
根据 $s$-$m$-$n$ 定理，存在一个函数 $s(e,x,y)$，使得对于每个 $z$，
\begin{align*}
\cfind{s(e,x,y)}(z) & \simeq \cfind{e}[3](x,y,z) \\
& \simeq \cfind{x}(y).
\end{align*}

\begin{explain}
根据上面的非形式论证，$s(e,x,y)$ 是如下机器的指标：对于任何输入 $z$，忽略该输入并计算 $\cfind{x}(y)$。
\end{explain}

特别地，我们有
\[
\cfind{s(e,x,y)}(0) \fdefined \quad \text{当且仅当} \quad
\cfind{x}(y) \fdefined.
\]
换言之，$\tuple{x, y} \in K_0$ 当且仅当 $s(e,x,y) \in
K_1$。所以，由
\[
g(w) = s(e,(w)_0,(w)_1)
\]
定义的函数 $g$ 是从 $K_0$ 到 $K_1$ 的一个归约。
\end{proof}

\end{document}